\scanpage{182}
d'un préschéma de Picard \ldots] \deleted{\unsure{Ceci montre}.}

Utilisant ce résultat, on voit que \uline{si} on peut trouver un faisceau \inserted{\uline{inv}} fondamental \deleted{\unsure{un} $\underline\Lambda$} $\underline\Lambda$ sur $\underline{\underline{\operatorname{Pic}}}\times_S X$ \deleted{Alors ce faisceau} [par exemple si $X/S$ a une section \ldots], alors \deleted{$\underline{\underline{\operatorname{Div}}}_{X/S}$ s'identifie à}
\[
\underline{\underline{\operatorname{Div}}}_{X/S}\simeq
\underline{\underline{\operatorname{Div}}}^{\underline\Lambda}_{\underline{\underline{\operatorname{Pic}}}_{X/S}\times_S X\,/\,\underline{\underline{\operatorname{Pic}}}_{X/S}}
\]

Cela prouve l'intérêt d'une étude directe, indépendamment de celle de $\underline{\underline{\operatorname{Pic}}}_{X/S}$, des foncteurs \inserted{de la forme} $\underline{\underline{\operatorname{Div}}}_{X/S}^{\underline L}$ et des \deleted{\unsure{leurs représentants|leurs représentant}} préschémas qui dans les cas favorables les représentent.

On va supposer maintenant que $X$ est propre \inserted{et plat} sur $S$, et que l'on a un isom fonctoriel :
\[
\underline H_f^0(\underline L\otimes_{\underline O_S}\underline M)=\underline{\operatorname{Hom}}_{O}(\underline E,\underline M)
\]
pour $\underline M$ \unsure{qu. coh.} sur $S$ [ce qui est vrai si $S$ est local, \unsure{ou|où|et} $X$ projectif sur $S$] \deleted{\unsure{$X$ quasi libre sur $S$} \ldots]}.

La connaissance du faisceau cohérent $\underline E$ sur $S$ \deleted{fixe}\inserted{détermine} le ``foncteur'' $S'\rightsquigarrow\underline H^0_{f'}(\underline L')$. \uline{Supposons p. ex.} que
% archive sheet number: 181
